\item \textbf{Test so sánh trước và sau thao tác.} Các giá trị như \texttt{ref} của console, pid hay \texttt{nsched} phụ thuộc vào \texttt{init}, \texttt{sh} và thời điểm chạy, nên mỗi test chụp \texttt{sysinfo} trước và sau thao tác rồi kiểm tra độ chênh lệch (ví dụ \texttt{nfile} tăng đúng 5 sau 5 lần \texttt{open}). Chỉ các hằng số (\texttt{NOFILE}, \texttt{NFILE}, \texttt{NPROC}, \texttt{init} có pid 1) mới được so sánh trực tiếp. Cuối mỗi test luôn kiểm tra \texttt{nfile}, \texttt{nproc} đã trở về như lúc đầu.

\item \textbf{Chỉ đối chiếu hai danh sách khi hệ thống đứng yên.} Do kết quả không phải bản chụp tức thời (ý đầu tiên của mục này), các kiểm tra đối chiếu danh sách file với danh sách process chỉ được làm khi mọi process khác đang bị chặn. Khi chờ một process đổi trạng thái, test chụp thêm một lần sau khi thấy trạng thái đó rồi mới đối chiếu.

\item \textbf{Giữ process con đứng yên bằng pipe.} Con bị chặn ở \texttt{read()} trên một pipe rỗng; cha ghi 1 byte khi muốn con chạy tiếp và thoát. Cách này chắc chắn hơn chờ bằng \texttt{sleep}. Khi cần chờ một trạng thái (\texttt{ZOMBIE}, \texttt{SLEEPING}), test gọi \texttt{sysinfo} liên tục tới khi thấy, tối đa 100 tick.

\item \textbf{Không phụ thuộc số CPU.} Test 10 chỉ yêu cầu các con chạy vòng lặp vô hạn luôn ở \texttt{RUNNING} hoặc \texttt{RUNNABLE} và có lúc thấy \texttt{RUNNABLE}; với 1 CPU thì CPU đó đang chạy chính \texttt{sysinfotest} nên không con nào \texttt{RUNNING}. Nhóm đã chạy với \texttt{CPUS=3} và \texttt{CPUS=1}, đều đạt.
